\section{Círculo}

Um círculo é dado pelo centro $o$ e pelo raio $r$. Todas as funções precisam de raiz quadrada, então \texttt{T} tem
que ser ponto flutuante (o \texttt{assert} garante isso). As interseções e as tangentes devolvem um \texttt{vector}
com os pontos encontrados.

\subsection*{Incentro e circuncentro}

\begin{drawing}
  \draw[fig_highlight] (2,1) circle (2.236);
  \fill[fig_fill] (1.460,1.052) circle (1.052);
  \draw[fig_accent] (1.460,1.052) circle (1.052);
  \draw[thick] (0,0) -- (4,0) -- (1,3) -- cycle;
  \fill (0,0) circle (1.5pt) node[below left] {$a$};
  \fill (4,0) circle (1.5pt) node[below right] {$b$};
  \fill (1,3) circle (1.5pt) node[above left] {$c$};
  \fill[fig_accent] (1.460,1.052) circle (2pt) node[left] {$i$};
  \draw[fig_accent, dashed] (1.460,1.052) -- (1.460,0);
  \fill[fig_highlight] (2,1) circle (2pt) node[right] {$o$};
  \draw[fig_highlight, dashed] (2,1) -- (4,0);
\end{drawing}

\texttt{incircle(a, b, c)} devolve o centro $i$ do círculo inscrito (azul), tangente aos três lados. Ele é a média
dos vértices com peso igual ao lado oposto: $i = (a \cdot |bc| + b \cdot |ac| + c \cdot |ab|) / (|ab| + |bc| + |ac|)$.
O raio é a distância de $i$ a qualquer lado, ou $2 \cdot \text{área} / \text{perímetro}$.

\texttt{circumcircle(a, b, c)} devolve o centro $o$ do círculo que passa pelos três vértices (vermelho), o encontro das
mediatrizes. Com $v = b - a$ e $w = c - a$,
\[ o = a + \frac{\text{perp}(v \, |w|^2 - w \, |v|^2)}{2 \, \text{cross}(v, w)}, \]
e o raio é \texttt{sqrt(dist2(o, a))}. As duas devolvem $(0, 0)$ se o triângulo é degenerado (pontos colineares).

\subsection*{Menor círculo que cobre os pontos}

\begin{drawing}
  \fill[fig_fill] (0,0) circle (2);
  \draw[fig_accent] (0,0) circle (2);
  \foreach \p in {(0.5,0.6), (-0.6,-0.2), (0.9,-0.3), (1.2,1.1), (-1,0.9), (0.1,-1.2)} \fill \p circle (1.5pt);
  \fill[fig_highlight] (-0.347,1.970) circle (2pt);
  \fill[fig_highlight] (-1.532,-1.286) circle (2pt);
  \fill[fig_highlight] (1.879,-0.684) circle (2pt);
  \draw[gray, dashed] (0,0) -- (1.879,-0.684);
  \node[gray] at (1.05,-0.15) {$r$};
  \fill[fig_accent] (0,0) circle (1.5pt) node[left] {$o$};
\end{drawing}

\texttt{min\_enclosing\_circle(p)} devolve \texttt{\{o, r\}}, o menor círculo com todos os pontos dentro. Ele é
definido por 2 ou 3 pontos da borda (vermelhos). O algoritmo de Welzl embaralha os pontos e vai aumentando o círculo:
se $p_i$ cai fora do círculo dos anteriores, $p_i$ está na borda do novo, então recomeça com o círculo só de $p_i$ e
adiciona os anteriores; se $p_j$ também cai fora, os dois estão na borda (círculo de diâmetro $p_i p_j$); se $p_k$ cai
fora, o círculo é o \texttt{circumcircle} dos três. Como a ordem é aleatória, a chance de $p_i$ cair fora é no máximo
$3/i$, e o total esperado é $O(n)$.

\subsection*{Círculo e reta}

\begin{drawing}
  \draw (0,0) circle (1.6);
  \fill (0,0) circle (1.5pt) node[below] {$o$};
  \draw[thick] (-2.6,0.8) -- (2.6,0.8) node[right] {2 pontos};
  \draw[thick] (-2.6,1.6) -- (2.6,1.6) node[right] {1 ponto};
  \draw[thick] (-2.6,2.3) -- (2.6,2.3) node[right] {nenhum};
  \draw[gray, dashed] (0,0) -- (0,0.8);
  \draw[gray] (0,0) -- (1.386,0.8);
  \node[gray] at (0.95,0.25) {$r$};
  \fill (0,0.8) circle (1.5pt) node[below left] {$p$};
  \draw[fig_accent, very thick] (0,0.8) -- (1.386,0.8);
  \node[fig_accent] at (0.69,1.08) {$h$};
  \fill[fig_highlight] (-1.386,0.8) circle (2pt);
  \fill[fig_highlight] (1.386,0.8) circle (2pt);
  \fill[fig_highlight] (0,1.6) circle (2pt);
\end{drawing}

\texttt{circle\_line\_inter(o, r, l)}: $p$ é a projeção de $o$ na reta e a distância de $o$ à reta é
$|\text{side}(o)| / |v|$. Por Pitágoras, $h^2 = r^2 - \text{dist}^2$. Se $h^2 < 0$ a reta passa por fora, se
$h^2 = 0$ ela é tangente e a resposta é só $p$, e se $h^2 > 0$ os pontos são $p \pm h$ na direção de $v$.

\subsection*{Círculo e segmento}

\begin{drawing}
  \draw (0,0) circle (1.6);
  \fill (0,0) circle (1.5pt) node[below] {$o$};
  \draw[gray, dashed] (-2.6,0.8) -- (3.4,0.8);
  \draw[thick] (-0.5,0.8) -- (2.6,0.8);
  \fill (-0.5,0.8) circle (1.5pt) node[above] {$a$};
  \fill (2.6,0.8) circle (1.5pt) node[above] {$b$};
  \draw[fig_highlight, fill=white] (-1.386,0.8) circle (2pt);
  \fill[fig_highlight] (1.386,0.8) circle (2pt);
\end{drawing}

\texttt{circle\_segment\_inter(o, r, a, b)} intersecta o círculo com a reta $ab$ e mantém só os pontos dentro do
segmento (\texttt{in\_disk}). O ponto vazado está na reta, mas fora de $[a, b]$. Se $a = b$, o segmento é um ponto e
a resposta é ele mesmo se estiver na circunferência.

\subsection*{Dois círculos}

\begin{drawing}
  \draw (0,0) circle (2);
  \draw (3,0) circle (1.6);
  \fill (0,0) circle (1.5pt) node[below left] {$o_1$};
  \fill (3,0) circle (1.5pt) node[below right] {$o_2$};
  \draw[gray, dashed] (0,0) -- (3,0);
  \draw[gray] (0,0) -- (1.74,0.986);
  \node[gray] at (0.7,0.65) {$r_1$};
  \draw[gray] (3,0) -- (1.74,0.986);
  \node[gray] at (2.6,0.75) {$r_2$};
  \fill (1.74,0) circle (1.5pt);
  \node at (2.2,-0.35) {$p$};
  \draw[fig_accent, very thick] (1.74,-0.986) -- (1.74,0.986);
  \node[fig_accent] at (1.25,0.3) {$h$};
  \fill[fig_highlight] (1.74,0.986) circle (2pt);
  \fill[fig_highlight] (1.74,-0.986) circle (2pt);
\end{drawing}

\texttt{circle\_inter(o1, r1, o2, r2)}: com $v = o_2 - o_1$ e $d = |v|$, os pontos de interseção ficam na
perpendicular a $v$ que passa por $p$. Igualando $h^2$ nos dois triângulos retângulos,
$|o_1 p| = (d^2 + r_1^2 - r_2^2) / 2d$ e $h^2 = r_1^2 - |o_1 p|^2$. O código guarda $\text{op} = d \cdot |o_1 p|$
para não dividir antes da hora. Como antes, $h^2 < 0$ não tem interseção, $h^2 = 0$ é tangente, e senão os pontos
são $p \pm h$ na direção de \texttt{perp(v)}. Círculos com o mesmo centro devolvem vazio (mesmo se forem iguais).

\subsection*{Área da interseção}

\begin{drawing}
  \begin{scope}
    \clip (0,0) circle (2);
    \fill[fig_fill] (3,0) circle (1.6);
  \end{scope}
  \draw (0,0) circle (2);
  \draw (3,0) circle (1.6);
  \fill (0,0) circle (1.5pt) node[left] {$o_1$};
  \fill (3,0) circle (1.5pt) node[right] {$o_2$};
  \draw[gray, dashed] (1.74,0.986) -- (0,0) -- (1.74,-0.986);
  \draw[gray, dashed] (1.74,0.986) -- (3,0) -- (1.74,-0.986);
  \draw[fig_accent] (1.74,0.986) -- (1.74,-0.986);
  \draw[fig_highlight] (0.6,-0.34) arc (-29.5:29.5:0.69);
  \node[fig_highlight] at (0.95,0) {$\theta_1$};
  \draw[fig_highlight] (2.645,0.277) arc (142:218:0.45);
  \node[fig_highlight] at (2.27,0) {$\theta_2$};
\end{drawing}

\texttt{circle\_inter\_area(o1, r1, o2, r2)}: a lente (azul) é a soma de dois segmentos circulares, um de cada
círculo, separados pela corda. O segmento do círculo $i$ com ângulo central $\theta_i$ tem área
$r_i^2 (\theta_i - \sin \theta_i) / 2$, e $\theta_1 = 2 \arccos\big((d^2 + r_1^2 - r_2^2) / 2 d r_1\big)$ vem da
lei dos cossenos (o \texttt{clamp} evita \texttt{acos} de $1 + \varepsilon$). Se $d \ge r_1 + r_2$ eles não se
tocam e a área é 0; se $d \le |r_1 - r_2|$ um está dentro do outro e a área é a do menor, $\pi \min(r_1, r_2)^2$.

\subsection*{Tangentes de um ponto}

\begin{drawing}
  \draw (0,0) circle (1.5);
  \fill (0,0) circle (1.5pt) node[left] {$o$};
  \fill (4,0) circle (1.5pt) node[right] {$p$};
  \draw[gray, dashed] (0,0) -- (4,0);
  \draw[thick] (4,0) -- (-0.141,1.675);
  \draw[thick] (4,0) -- (-0.141,-1.675);
  \draw[gray] (0,0) -- (0.5625,1.3905);
  \node[gray] at (0.1,0.85) {$r$};
  \draw[gray] (0.5625,0) -- (0.5625,1.3905);
  \fill (0.5625,0) circle (1.5pt) node[below right] {$q$};
  \fill[fig_highlight] (0.5625,1.3905) circle (2pt) node[above right] {$t_2$};
  \fill[fig_highlight] (0.5625,-1.3905) circle (2pt) node[below right] {$t_1$};
\end{drawing}

\texttt{circle\_tangents(o, r, p)} devolve os pontos onde as retas que saem de $p$ tocam o círculo. O triângulo
$o\,t\,p$ tem ângulo reto em $t$, então, com $v = p - o$ e $d = |v|$, a sombra de $t - o$ em $v$ mede $r^2 / d$
(o ponto $q$) e a altura mede $r \sqrt{d^2 - r^2} / d$, na direção de \texttt{perp(v)}. Se $p$ está dentro do
círculo não há tangente; se está na borda, a resposta é o próprio $p$.

\subsection*{Tangentes comuns}

\begin{drawing}
  \draw (0,0) circle (1.2);
  \draw (4.5,0) circle (0.8);
  \fill (0,0) circle (1.5pt) node[left] {$o_1$};
  \fill (4.5,0) circle (1.5pt) node[right] {$o_2$};
  \draw[fig_accent, thick] (-0.563,1.255) -- (5.241,0.737);
  \draw[fig_accent, thick] (-0.563,-1.255) -- (5.241,-0.737);
  \draw[fig_highlight, thick] (-0.008,1.344) -- (4.686,-0.985);
  \draw[fig_highlight, thick] (-0.008,-1.344) -- (4.686,0.985);
  \foreach \p in {(0.107,1.195), (4.571,0.797), (0.107,-1.195), (4.571,-0.797)} \fill[fig_accent] \p circle (2pt);
  \foreach \p in {(0.533,1.075), (4.144,-0.717), (0.533,-1.075), (4.144,0.717)} \fill[fig_highlight] \p circle (2pt);
\end{drawing}

\texttt{circle\_circle\_tangents(o1, r1, o2, r2)} devolve cada tangente como o par de pontos onde ela toca os dois
círculos, então \texttt{line(a, b)} é a reta. As externas (azul) deixam os círculos do mesmo lado e as internas
(vermelho) passam entre eles. Se $w$ é a normal unitária da reta, a distância dela aos centros dá
\texttt{dot(w, v)} $= r_1 - r_2$ nas externas; as internas são a mesma conta trocando $r_2$ por $-r_2$, e o laço em
\texttt{s} faz as duas. Separados há 4, tangentes por fora 3, se cortando 2, tangentes por dentro 1, um dentro do
outro 0. Quando os círculos se tocam, os dois pontos do par são iguais e a reta é \texttt{line(a, a + perp(v))}.
